% Segmentation Pipeline section. Paste into your own thesis, or \input it.

\section{Data and preprocessing}
\label{sec:pipeline-data}

The network's input is a CCTA volume; its target is a binary \textbf{lumen mask} -- a 3D image
where each tiny volume element (a voxel, the 3D equivalent of a pixel) is labelled as either
``blood'' or ``not blood,'' marking out the hollow, blood-filled channel of the coronary arteries. The scan comes from the original
ImageCAS cohort \cite{zeng2023imagecas}; the mask is ImageCAS-X's own re-annotation of that
cohort \cite{bransby2026imagecasx}, which labels 14 individual coronary segments separately.
Since segmentation here only needs to find the vessel, not identify which segment it belongs to,
those 14 labels are collapsed to a single foreground label (vessel vs.\ background) before
training.

Every scan and mask is first resampled, once offline, onto a common \textbf{0.5\,mm isotropic}
grid. Voxel spacing is the 3D equivalent of a photo's pixel size: it says how much real-world
distance one voxel covers, the same way a camera's resolution says how much of a scene one pixel
covers. Not every scan in the cohort was taken at the same spacing, so without this step the same
stretch of artery would be built out of a different number of voxels from one scan to the next.

The cohort's 800 usable cases -- 200 of the original 1000 are excluded for low image quality
(Figure~\ref{fig:pipeline-imagecasx}) -- are split once into disjoint \textbf{train},
\textbf{validation} and
\textbf{test} sets. Training updates the network's weights; validation is watched each epoch to
pick which checkpoint to keep; test is touched only once, at the end, to produce the numbers
reported in Table~\ref{tab:pipeline-benchmark}.

\begin{figure}[htbp]
  \centering
  \includegraphics[width=\linewidth=0.5]{content/background/figures/graphical_abstract.png}
  \caption{ImageCAS-X at a glance \cite{bransby2026imagecasx}. Of the 1000 ImageCAS CCTA volumes,
    800 are kept and 200 excluded for image-quality problems (cardiac motion, step artefacts,
    sub-voxel contrast). The kept scans are re-annotated with a centerline tree, a lumen mask
    labelled by coronary segment, a surface mesh and scan descriptors, and scored in the
    segmentation benchmark on the right. Graphical abstract from the ImageCAS-X repository.}
  \label{fig:pipeline-imagecasx}
\end{figure}

The network never sees a whole scan during training, only small $128\times160\times160$ crops --
a full volume does not fit a GPU's memory alongside a 3D network. Since a coronary artery occupies
a small fraction of a volume, crop origins are not sampled uniformly at random: a configurable
fraction of each batch (0.5 for CAS-Net) is forced onto a location containing GT vessel, so the
network still sees enough foreground. Each scan's cached array is itself a \textbf{stack of 2D axial slices} -- one thin cross-sectional
image per position along the body's head-to-foot axis -- concatenated into a single 3D block. How
many slices a scan contains is not fixed across the cohort: it follows however many axial images
that subject's original CCTA acquisition produced (206, for one representative case), so stack
length varies from one subject to the next even though every scan now shares the same 0.5\,mm
spacing after resampling. The $128\times160\times160$ crop is a fixed-size sub-block cut out of
that variable-length stack at a randomly chosen, foreground-biased position (160 slices deep); if
a scan happens to be smaller than the crop along some axis, that side is zero-padded rather than
dropped, so the network always receives a crop of exactly the configured size regardless of the
source scan's length.



\begin{figure}[htbp]
  \centering
  \includegraphics[width=\linewidth=0.5]{content/background/figures/val_batch0.png}
  \caption{Dataloader sanity check on one training batch of five crops. Top row: the single axial
    slice sitting at the middle of each crop's 160-slice stack. Bottom row: the same slice with the
    ground-truth (GT) lumen mask overlaid as a boundary outline rather than a filled region, so the
    CT tissue underneath stays visible. Small yellow rings mark where the GT says a coronary artery
    cross-section sits in that slice.}
  \label{fig:pipeline-dataloader-check}
\end{figure}



\textbf{An overview just for me}

% WORKING NOTE -- \ref above and this \input are the two places to touch when removing
% the scratch table before final submission; see 01a_data_storage_notes.tex.
% WORKING NOTE -- not part of the thesis narrative, kept only for our own reference
% while the pipeline is being built out. Plug out before final submission by deleting
% the \input line in 01_data_preparation.tex that pulls this file in (and, optionally,
% this file itself). Needs tabularx, booktabs -- inherited from the including document.

\begingroup
\footnotesize
\renewcommand{\arraystretch}{0.85}

\begin{table}[htbp]
  \centering
  \begin{tabularx}{\textwidth}{@{}l X@{}}
    \toprule
    \textbf{Data} & \textbf{Where \& how stored} \\
    \midrule
    Raw CCTA volumes &
      \texttt{imagecasx\_data/volumes/}, local copy, 1000 scans \\
    Original ImageCAS labels &
      \texttt{imagecasx\_data/imagecas\_orig\_labels/}, local copy, unused, 1000 scans \\
    ImageCAS-X re-annotation &
      \texttt{imagecasx\_data/segmentations/} (+ centerlines/surfaces/filelist/
      Descriptors.xlsx), symlinked from \texttt{ImageCAS-X\_dataset/}, 800 scans \\
    \bottomrule
  \end{tabularx}
  \caption{Working note: how the delivered data is stored. Not for the final thesis.}
  \label{tab:data-storage-notes}
\end{table}

\begin{table}[htbp]
  \centering
  \begin{tabularx}{\textwidth}{@{}X X@{}}
    \toprule
    \textbf{Command} & \textbf{Produces / reads} \\
    \midrule
    \texttt{python -m utils.offline\_resample\_\allowbreak images\_to\_disk -c
      configs/cas\_net.json --workers 8} &
      writes \texttt{volumes\_resampled/}, \texttt{segmentations\_resampled/}
      (\texttt{.npy} cache, 800 scans) \\
    \texttt{bsub < jobs/train\_cas\_net.sh} &
      reads \texttt{\$ImageCAS\_X\_data\_path} (\texttt{env.sh}); override with
      \texttt{-env "all, CONFIG=\allowbreak configs/cas\_net\_smoke.json"} for the
      5-epoch smoke run \\
    \bottomrule
  \end{tabularx}
  \caption{Working note: commands used in the CAS-Net data pipeline. Not for the
    final thesis.}
  \label{tab:data-commands-notes}
\end{table}

\endgroup


A CT scan measures density on a fixed brightness scale (Hounsfield units, HU), where air is very
dark, soft tissue and contrast-filled vessel sit in the middle, and bone or metal are very bright.
Only that middle range actually helps tell vessel from background, so on every patch load the
scan undergoes \textbf{intensity normalisation}: its values are cut off at a fixed range ($-200$ to
$1000$\,HU) -- throwing away the very darkest and very brightest extremes -- and the rest is
rescaled to run from 0 to 1.

During training only, each crop is further passed through a chain of random augmentations before
reaching the network, detailed in Section~\ref{sec:pipeline-augmentation} alongside the rest of
the training loop.
